% ============================================================================
% Physical AI Oncology Clinical Trial Site Documentation
% Combined LaTeX Source File - All 11 Documents
% Version: 2.9.0 | March 24, 2026
% Author: Kevin Kawchak, CEO, ChemicalQDevice (kevink@chemicalqdevice.com)
% Adapted using Claude Code Opus 4.6
% ============================================================================
%
% This file contains the complete LaTeX source for all 11 site documents.
% Each document section is clearly delimited. To compile individual documents,
% extract the relevant section and use it with its corresponding .bib and .sty
% files from the document's subdirectory.
%
% Document Index:
%   1. SB 1042 - California Physical AI Trial Authorization Act
%   2. AB 2847 - California Physical AI Patient Rights and Robotic Safety Act
%   3. SB 892 - California Physical AI Clinical Data Protection Act
%   4. San Francisco Municipal Code Update
%   5. California Code of Regulations Title 22
%   6. FDA National Compliance Guide
%   7. Building Code Standards
%   8. Premises Code
%   9. Parking and Transportation Standards
%  10. Site Activation and Standard Operating Procedures
%  11. Emergency Preparedness Plan
%
% Disclaimers:
%   - Unofficial independent draft derived from prior legislation; no third-party
%     endorsement, sponsorship, affiliation, or authorization is expressed or
%     implied; and was adapted using Claude Code Opus 4.6.
%   - This work was adapted from original CFR documents in the public domain.
%     The original ICH document is copyrighted and may be used, reproduced,
%     incorporated into other works, adapted, modified, translated or distributed
%     under a public license. This current work is not endorsed or sponsored by
%     CFR, ICH, or FDA; and was adapted using Claude Code Opus 4.6.
%   - This draft is independent and is not endorsed, sponsored, or approved by
%     any trial sponsor, CRO, site, IRB, regulator, or medical society; and was
%     adapted using Claude Code Opus 4.6.
%
% ============================================================================

\documentclass[11pt,letterpaper]{article}

% --- Shared Packages ---
\usepackage[margin=1in]{geometry}
\usepackage{setspace}
\onehalfspacing
\usepackage[T1]{fontenc}
\usepackage{lmodern}
\usepackage{microtype}
\usepackage{titlesec}
\titleformat{\section}{\large\bfseries}{\thesection.}{0.5em}{}
\titleformat{\subsection}{\normalsize\bfseries}{\thesubsection.}{0.5em}{}
\usepackage{enumitem}
\setlist[enumerate]{itemsep=2pt,parsep=2pt}
\setlist[itemize]{itemsep=2pt,parsep=2pt}
\usepackage{fancyhdr}
\pagestyle{fancy}
\fancyhf{}
\fancyhead[L]{\small Physical AI Oncology Trial Site Documentation}
\fancyhead[R]{\small March 2026}
\fancyfoot[C]{\thepage}
\renewcommand{\headrulewidth}{0.4pt}
\renewcommand{\footrulewidth}{0pt}
\usepackage[colorlinks=true,linkcolor=blue!60!black,citecolor=blue!60!black,urlcolor=blue!60!black]{hyperref}
\usepackage{tocloft}
\setlength{\cftsecnumwidth}{2.5em}
\setlength{\cftsubsecnumwidth}{3.2em}
\usepackage{booktabs}
\usepackage{longtable}
\usepackage{xcolor}
\tolerance=200
\emergencystretch=2em
\hyphenpenalty=50
\widowpenalty=10000
\clubpenalty=10000

\title{%
  \textbf{Physical AI Oncology Clinical Trial Site}\\
  \textbf{Complete Documentation Package}\\[6pt]
  \large 11 Documents for California's First Physical AI\\
  Oncology Clinical Trial Site\\[6pt]
  \normalsize Legislation, Regulations, Building Code,\\
  Premises Code, and Operations
}
\author{%
  Kevin Kawchak\\
  CEO, ChemicalQDevice\\
  \texttt{kevink@chemicalqdevice.com}
}
\date{March 24, 2026}

\begin{document}
\maketitle
\thispagestyle{fancy}

\begin{center}
\footnotesize
\textit{Unofficial independent draft derived from prior legislation;
no third-party endorsement, sponsorship, affiliation, or authorization
is expressed or implied; and was adapted using Claude Code Opus 4.6.}
\end{center}

\vspace{0.5em}

\begin{center}
\footnotesize
\textit{This work was adapted from original CFR documents in the public domain.
The original ICH document is copyrighted and may be used, reproduced,
incorporated into other works, adapted, modified, translated or distributed
under a public license. This current work is not endorsed or sponsored by
CFR, ICH, or FDA; and was adapted using Claude Code Opus 4.6.}
\end{center}

\vspace{0.5em}

\begin{center}
\footnotesize
\textit{This draft is independent and is not endorsed, sponsored, or approved
by any trial sponsor, CRO, site, IRB, regulator, or medical society;
and was adapted using Claude Code Opus 4.6.}
\end{center}

\vspace{1em}
\tableofcontents
\newpage

% ============================================================================
% DOCUMENT 1: SB 1042 - Trial Authorization Act
% Source: 01-legislation-authorization/physical_ai_trial_authorization.tex
% ============================================================================

\part{SB 1042 - California Physical AI Oncology Clinical Trial
Authorization and Site Establishment Act of 2026}

\section{Legislative Findings and Declarations}

The Legislature finds and declares all of the following:

\begin{enumerate}[label=(\alph*)]
\item Physical artificial intelligence represents a transformative advancement
  in oncology clinical care delivery, including robotic surgical systems,
  collaborative robots, radiotherapy positioning systems, imaging assistants,
  rehabilitation exoskeletons, and companion robots.
\item A 24-hour simulation demonstrated 168 unique patients across 15 cancer
  types served by 29 robot instances in 24 hours (10x traditional throughput).
\item The simulation achieved PSL 63.4-64.4 (Advanced Site), 99.7\% uptime,
  8-minute average wait time, and 4.7/5.0 patient satisfaction.
\item All 7 adverse events were Grade 1-2, resolved same hour, zero patient
  harm. Both minute-level and 24-hour details handled correctly by AI.
\item The format shifts control toward the patient: 24/7 scheduling, 75-85\%
  cost reduction, 85-95\% retention (vs 70-80\% traditional).
\item Staffing reduced from 80-120 FTE to 8-10 FTE; more patients treated with
  more robots and fewer workers.
\end{enumerate}

\section{Definitions}

Key definitions include Physical AI system (10 categories), PSL (three
regulatory dimensions: Omniscient, Omnipresent, Omnipotent), USL (four
dimensions: simulation switching, AI integration, cross-robot sharing,
clinical trial collaboration), on-demand trial format, digital twin, MCP-PAI
standard, qualified site, and sim-to-real validation.

\section{Authorization Process}

CDPH shall establish a process for authorizing Physical AI oncology clinical
trial sites in California, beginning with San Francisco. Applications include
facility plans, equipment specifications, staffing plans, USL scores, Phase 0
simulation validation, regulatory compliance plans, and emergency response
plans. The department shall act within 120 days.

\section{Safety, Reporting, and Workforce}

The act establishes adverse event reporting (7-day and 15-day timelines),
Physical AI Safety Review Committee, hash-chained audit trails, and workforce
transition programs. It relates to existing California AI legislation: AB 489
(2025), AB 3030 (2024), SB 1120 (2024), SB 243 (2025), AB 2013 (2024).

\newpage

% ============================================================================
% DOCUMENT 2: AB 2847 - Patient Rights and Robotic Safety Act
% Source: 02-legislation-patient-rights/physical_ai_patient_rights.tex
% ============================================================================

\part{AB 2847 - California Physical AI Patient Rights
and Robotic Safety Act of 2026}

\section{Patient Bill of Rights}

Six categories of protected rights for Physical AI oncology trial patients:

\begin{enumerate}[label=(\arabic*)]
\item \textbf{Right to On-Demand Scheduling:} 24/7 procedure scheduling with
  maximum 30-minute wait. Simulation demonstrated 8-minute average.
\item \textbf{Right to Full Disclosure:} Robot type, model, role, autonomy mode,
  USL score in plain language, safety record, emergency stop provisions,
  digital twin usage, and human oversight qualifications.
\item \textbf{Right to Human Intervention:} Request human intervention within 2
  minutes at any point during a procedure. Withdraw consent for Physical AI.
\item \textbf{Right to Privacy and Data Control:} HIPAA, CCPA, CPRA compliance.
  Deny-by-default MCP-PAI authorization. Right to digital twin deletion.
\item \textbf{Right to Equitable Access:} No discrimination. Accommodations for
  non-traditional hours, caregiving, transportation barriers.
\item \textbf{Pediatric Protections:} Age-appropriate demonstrations, parental
  consent, companion robot protocols, advocate appointments for wards.
\end{enumerate}

\section{Robotic Safety Standards}

Physical AI safety certification before deployment, continuous telemetry
monitoring, force limit enforcement, and same-hour adverse event resolution
targeting the simulation's 100\% same-hour resolution rate.

\newpage

% ============================================================================
% DOCUMENT 3: SB 892 - Data Protection and Transparency Act
% Source: 03-legislation-data-transparency/physical_ai_data_transparency.tex
% ============================================================================

\part{SB 892 - California Physical AI Clinical Data Protection
and Transparency Act of 2026}

\section{Data Collection and Protection}

Comprehensive notice requirements for all Physical AI clinical data categories.
Real-time patient data access through secure portal. AES-256 encryption at
rest, TLS 1.3 in transit. Deny-by-default authorization. Six MCP-PAI actor
roles. Hash-chained audit trails (SHA-256) with 15-year retention. HIPAA Safe
Harbor de-identification for all shared data.

\section{Federated Learning and AI Transparency}

Differential privacy with epsilon not exceeding 1.0 per round (validated in
70-round, 12-site cancer patient journey). Membership inference testing for
model parameters. PCCP filing for AI model updates. Five AI categories with
transparency requirements per ICH E6(R3).

\section{Enforcement}

Civil penalties: \$2,500-\$7,500 per violation per patient. Aggravated
violations (audit trail failure, reporting failures): up to \$25,000.
Attorney General enforcement authority.

\newpage

% ============================================================================
% DOCUMENT 4: San Francisco Municipal Code
% Source: 04-city-regulations/sf_city_regulations.tex
% ============================================================================

\part{San Francisco Municipal Code Update - Physical AI Trial Sites}

\section{Zoning and Permits}

Conditionally permitted in NC, MUG, PDR, DTG, C-3, and Medical Use districts.
Preference for locations within 1/4 mile of BART/Muni. Minimum 40,000 sq ft
lot. Maximum 65 feet height. Conditional use permit requires Physical AI System
Deployment Plan, Patient Flow Analysis, Neighborhood Impact Assessment,
Community Benefits Plan, and TDM Plan. Public notice to 300-foot radius.

\section{Building and Health Code}

HEPA filtration ISO Class 7 for surgical suites, 150 PSF floor loads, shielded
RT vaults. Health Care Facility Operating License specifying authorized robot
types and maximum patients. Quarterly inspections published within 30 days.
Noise limits 45 dBA nighttime at property line.

\newpage

% ============================================================================
% DOCUMENT 5: California State Regulations
% Source: 05-state-regulations/ca_state_regulations.tex
% ============================================================================

\part{California Code of Regulations, Title 22 - Physical AI Trial Sites}

\section{Authorization and Registration}

CDPH authorization process with 120-day timeline. \$25,000 filing fee.
3-year certificate with annual verification. Physical AI system registration
with USL scores. Minimum USL thresholds: 7.0 surgical, 6.0 therapeutic
positioning and diagnostic, 5.0 cobots, 4.0 rehabilitative and humanoid,
3.0 companion.

\section{Staffing, Training, and Enforcement}

Minimum 8-10 FTE across three shifts with on-call physician and pharmacist.
160-hour site safety officer training and CDPH certification. Quarterly
unannounced inspections. Enforcement: deficiency notices, \$10,000/day
penalties, system suspension, site authorization revocation.

\newpage

% ============================================================================
% DOCUMENT 6: FDA National Compliance Guide
% Source: 06-national-regulations/fda_national_regulations.tex
% ============================================================================

\part{FDA Physical AI Oncology Trial Site National Compliance Guide}

\section{Regulatory Correction Map}

Section-by-section correction guide for federal regulations:

\begin{itemize}
\item \textbf{21 CFR Part 50:} 5 sections (scope, definitions, informed consent,
  consent elements, pediatric protections).
\item \textbf{21 CFR Part 312:} 10 sections (scope, definitions, phases, IND
  content, amendments, safety reporting, annual reports, withdrawal,
  readiness, clinical holds).
\item \textbf{ICH E6(R3):} 7 areas (principles, investigator, sponsor, data,
  system documentation, protocol templates, glossary).
\item \textbf{21 CFR Part 11:} 3 sections (closed systems, open systems,
  signatures).
\item \textbf{21 CFR Part 820:} 4 sections (design, production, nonconforming,
  complaints).
\item \textbf{HIPAA:} 4 Physical AI-specific considerations.
\item \textbf{NRC 10 CFR Part 35:} 3 radiation safety considerations.
\end{itemize}

Cross-reference matrix maps all 10 robot types to applicable regulations.

\newpage

% ============================================================================
% DOCUMENT 7: Building Code Standards
% Source: 07-building-code/physical_ai_building_code.tex
% ============================================================================

\part{Physical AI Trial Facility Building Code Standards}

\section{Structural and Mechanical}

85,000-100,000 sq ft facility with 9 clinical wings supporting 29 robot
instances. Floor loads: 150 PSF surgical, 200 PSF RT vaults, 250 PSF servers.
Vibration: VC-E surgical, VC-D imaging, VC-C radiotherapy. HEPA ISO Class 7
surgical suites, 20 ACH, positive pressure. Server room cooling: 200 kW
rejection, N+1 redundancy, dedicated HVAC zone.

\section{Electrical and Technology}

Dual utility feeds, 800-1,200 kW. 72-hour generators. UPS: 30-min surgical,
15-min servers. Dedicated panels per wing with 25\% spare capacity. 10 Gbps
fiber backbone, segmented networks. Robot charging and docking infrastructure.
Clean agent fire suppression (FM-200/Novec 1230).

\newpage

% ============================================================================
% DOCUMENT 8: Premises Code
% Source: 08-premises-code/physical_ai_premises_code.tex
% ============================================================================

\part{Physical AI Trial Site Premises Code}

\section{Security and Access}

Five access control zones: Public, Patient, Clinical, Technical, Administrative.
24-hour video surveillance, electronic access control, intrusion detection.
Cybersecurity: robot control networks accessible only from Technical Zone, no
wireless access to control systems, MCP-PAI servers in locked facilities.

\section{Patient Flow and Robot Zones}

24-hour patient entrance. Covered drop-off zone for 3 vehicles. Automated
check-in. 20-seat reception area. 8-foot minimum corridor width. Separate
robot movement corridors with proximity sensors. Maintenance bay for 3
simultaneous systems. Emergency evacuation with automated robot safe-shutdown.

\newpage

% ============================================================================
% DOCUMENT 9: Parking and Transportation
% Source: 09-parking-transportation/physical_ai_parking.tex
% ============================================================================

\part{Physical AI Trial Site Parking and Transportation}

\section{Parking Facility}

120 spaces: 80 patient/visitor, 15 staff, 10 accessible, 5 drop-off, 10
ride-share. Maximum 3-level structure. Covered connection to building (300 ft
max). 24-hour monitoring with emergency call stations. 20\% EV Level 2 charging
(24 spaces), 4 DC fast charging, 20\% EV-capable conduit.

\section{Transportation}

Transit integration within 1/4 mile BART/Muni. Real-time transit displays.
Ride-share staging with geofenced pickup. 30 bicycle spaces (20 Class I, 10
Class II) with e-bike charging. Dedicated ambulance bay for 2 vehicles.
Sustainability: stormwater management, 50\% solar/green roof, LED with
occupancy sensors.

\newpage

% ============================================================================
% DOCUMENT 10: Site Operations
% Source: 10-site-operations/physical_ai_site_operations.tex
% ============================================================================

\part{Physical AI Trial Site Activation and SOPs}

\section{Activation Checklist}

Four categories: regulatory approvals (CDPH certificate, CUP, IRB, IND),
facility readiness (systems commissioned, MCP-PAI configured, power/HVAC
tested), staffing readiness (officers certified, on-call agreements), and
documentation (SOPs, patient instructions in multiple languages, consent
forms, emergency plan).

\section{Daily Operations}

Three shifts (06:00-14:00, 14:00-22:00, 22:00-06:00). Structured shift
handoff. Daily system verification checks (calibration, software integrity,
network, emergency stop, sensors). Utilization management based on simulation
baselines. Quality indicators tracked against 168-patient, 99.7\%-uptime,
8-minute wait baseline. Same-hour adverse event resolution target.

\newpage

% ============================================================================
% DOCUMENT 11: Emergency Preparedness
% Source: 11-emergency-preparedness/physical_ai_emergency.tex
% ============================================================================

\part{Physical AI Trial Site Emergency Preparedness}

\section{Emergency Classification}

Four levels: Level 1 (Minor, single system/patient), Level 2 (Moderate,
multiple systems/patients, Grade 2-3 AE), Level 3 (Major, partial/complete
evacuation, Grade 4+ AE), Level 4 (Catastrophic, facility non-operational).

\section{Response Procedures}

Emergency stop protocol: cease motion, retract instruments, lower couches,
de-energize radiation, maintain monitoring, generate audit record.
Cybersecurity: 5 threat categories (control compromise, model poisoning,
digital twin manipulation, MCP-PAI attack, network breach). Seismic: automated
safe-mode at 0.1g. Power: dual feeds, UPS bridge, 72-hour generators.

\section{Business Continuity}

Mutual aid with 2+ hospitals within 30 minutes. Remote backup: RPO 15 min,
RTO 4 hours. Post-emergency recalibration and Phase 0 re-validation if needed.
Quarterly drills rotating through system failure, cybersecurity, evacuation,
and power failure scenarios.

\end{document}
